% ==============================================================================
% TSL-51: Thai Sign Language Recognition System
% ระบบจดจำภาษามือไทย TSL-51
% รายงานวิจัยฉบับสมบูรณ์ - 5 บท
% ==============================================================================

\documentclass[12pt,a4paper]{report}
\usepackage[margin=1in]{geometry}
\usepackage[utf8]{inputenc}
%\usepackage[Thai]{babel}
\usepackage{graphicx}
\usepackage{hyperref}
\usepackage{amsmath}
\usepackage{amssymb}
\usepackage{algorithm}
\usepackage[noend]{algpseudocode}
\usepackage{listings}
\usepackage{xcolor}
\usepackage{longtable}
\usepackage{booktabs}
\usepackage{caption}

% Simple setup - no fontspec for pdflatex compatibility

% Color definitions
\definecolor{codegreen}{rgb}{0,0.6,0}
\definecolor{codegray}{rgb}{0.5,0.5,0.5}
\definecolor{codepurple}{rgb}{0.58,0,0.82}
\definecolor{backcolour}{rgb}{0.95,0.95,0.92}
\definecolor{navy}{RGB}{0,0,128}

\hypersetup{
    colorlinks=true,
    linkcolor=navy,
    filecolor=magenta,
    urlcolor=navy,
}

% Code listing style
\lstdefinestyle{mystyle}{
    backgroundcolor=\color{backcolour},
    commentstyle=\color{codegreen},
    keywordstyle=\color{blue},
    numberstyle=\tiny\color{codegray},
    stringstyle=\color{codepurple},
    basicstyle=\ttfamily\small,
    breaklines=true,
    frame=single,
    numbers=left
}
\lstset{style=mystyle}

% ==============================================================================
% DOCUMENT SETTINGS
% ==============================================================================
\title{
    \textbf{TSL-51: ระบบจดจำภาษามือไทย} \\
    \textit{Thai Sign Language Recognition System using Deep Learning}
}
\author{
    \textbf{TSL Research Team} \\
    \textit{ทีมวิจัยระบบจดจำภาษามือไทย}
}
\date{\textbf{เมษายน 2569}}

% ==============================================================================
% MAIN DOCUMENT
% ==============================================================================
\begin{document}

\pagenumbering{roman}
\thispagestyle{empty}

\begin{center}
\vspace*{2in}
{\Huge\bfseries TSL-51: ระบบจดจำภาษามือไทย\par}
\vspace{0.5in}
{\Large\bfseries Thai Sign Language Recognition System\par}
\vspace{1in}
{\LARGE\bfseries รายงานวิจัย\par}
\vspace{0.5in}
{\Large บทที่ 1-5\par}
\vspace{1.5in}
{\Large\bfseries TSL Research Team\par}
{\Large ทีมวิจัยระบบจดจำภาษามือไทย\par}
\vspace{0.5in}
{\Large เมษายน 2569\par}
\end{center}

\clearpage

% ==============================================================================
% TABLE OF CONTENTS
% ==============================================================================
\tableofcontents
\clearpage

% ==============================================================================
% CHAPTER 1: บทนำ
% ==============================================================================
\pagenumbering{arabic}
\chapter{บทนำ}

\section{ความเป็นมาและความสำคัญของงานวิจัย}

ภาษามือ (Sign Language) เป็นรูปแบบการสื่อสารหลักสำหรับผู้บกพร่องทางการได้ยินทั่วโลก โดยแต่ละประเทศจะมีภาษามือที่แตกต่างกัน 
ในประเทศไทยมีผู้ใช้ภาษามือไทย (Thai Sign Language: TSL) จำนวนมาก ซึ่งมีลักษณะเฉพาะและแตกต่างจากภาษามือในประเทศอื่น 
การแปลภาษามือเป็นข้อความอัตโนมัติ (Automatic Sign Language Translation) จึงมีความสำคัญอย่างยิ่งในการช่วยให้ผู้บกพร่องทางการได้ยินสามารถสื่อสารกับผู้อื่นได้สะดวกยิ่งขึ้น

ปัจจุบันเทคโนโลยีปัญญาประดิษฐ์ (Artificial Intelligence: AI) โดยเฉพาะอย่างยิ่ง Deep Learning ได้พัฒนาไปอย่างก้าวกระโดด 
ทำให้การจดจำท่าทาง (Gesture Recognition) และการจดจำภาษามือ (Sign Language Recognition) มีความแม่นยำมากขึ้น 
อย่างไรก็ตาม งานวิจัยในด้านภาษามือไทยยังมีจำนวนน้อยเมื่อเทียบกับภาษามือของประเทศตะวันตก เช่น American Sign Language (ASL) 
หรือ British Sign Language (BSL)

งานวิจัยนี้นำเสนอระบบ TSL-51 (Thai Sign Language 51 Signs) สำหรับการจดจำท่าทางมือในภาษามือไทยจำนวน 51 ท่าทาง 
ซึ่งครอบคลุมท่าทางพื้นฐานที่ใช้ในชีวิตประจำวัน ระบบที่พัฒนาขึ้นสามารถทำงานได้ทั้งในโหมดออฟไลน์ (Offline) 
และแบบเรียลไทม์ (Real-time) โดยใช้เทคนิค Deep Learning ร่วมกับ MediaPipe สำหรับการสกัดคุณลักษณะจากวิดีโอ

\section{วัตถุประสงค์ของงานวิจัย}

วัตถุประสงค์หลักของงานวิจัยนี้คือ:

\begin{enumerate}
    \item \textbf{พัฒนาระบบจดจำภาษามือไทย TSL-51} ที่สามารถจดจำท่าทางมือได้อย่างถูกต้องและรวดเร็ว
    \item \textbf{สร้างส่วนติดต่อผู้ใช้แบบ TUI} (Text User Interface) สำหรับการใช้งานระบบอย่างสะดวก
    \item \textbf{พัฒนาระบบแปลวิดีโอเป็นข้อความแบบเรียลไทม์} โดยใช้กล้องเว็บแคม
    \item \textbf{ทดสอบและเปรียบเทียบประสิทธิภาพ} ของโมเดลต่างๆ ได้แก่ GRU และ MLP
    \item \textbf{เปิดเผยชุดข้อมูล TSL-51} ให้เป็นสาธารณะบน HuggingFace สำหรับการวิจัยต่อยอด
\end{enumerate}

\section{ขอบเขตของงานวิจัย}

งานวิจัยนี้มีขอบเขตดังนี้:

\begin{itemize}
    \item \textbf{จำนวนท่าทาง:} 51 ท่าทางพื้นฐานของภาษามือไทย
    \item \textbf{ชุดข้อมูล:} TSL-51 Dataset จาก HuggingFace (547 วิดีโอต้นฉบับ)
    \item \textbf{โมเดลที่ใช้:} 
    \begin{itemize}
        \item GRU (Gated Recurrent Unit)
        \item MLP (Multi-Layer Perceptron)
    \end{itemize}
    \item \textbf{การประมวลผลวิดีโอ:} ใช้ MediaPipe สำหรับสกัดจุด Landmark บนมือและท่าทาง
    \item \textbf{จำนวน Feature:} 162 มิติ ประกอบด้วย:
    \begin{itemize}
        \item 63 มิติ สำหรับมือซ้าย (x, y, z ของ 21 จุด)
        \item 63 มิติ สำหรับมือขวา (x, y, z ของ 21 จุด)
        \item 36 มิติ สำหรับท่าทางร่างกาย (Pose landmarks)
    \end{itemize}
    \item \textbf{รูปแบบการทำงาน:}
    \begin{itemize}
        \item การ training แบบ batch
        \item การ inference แบบ single sample
        \item การ predict จากวิดีโอ
        \item การ translate แบบเรียลไทม์จากกล้อง
    \end{itemize}
\end{itemize}

\section{ความสำคัญของงานวิจัย}

ความสำคัญของงานวิจัยนี้มีดังนี้:

\begin{enumerate}
    \item \textbf{การช่วยเหลือผู้บกพร่องทางการได้ยิน:} ระบบที่พัฒนาขึ้นช่วยให้ผู้บกพร่องทางการได้ยินสื่อสารกับผู้อื่นได้สะดวกยิ่งขึ้น 
    โดยการแปลภาษามือเป็นข้อความแบบเรียลไทม์
    
    \item \textbf{การเรียนรู้ภาษามือ:} ระบบสามารถใช้เป็นเครื่องมือช่วยในการเรียนรู้ภาษามือไทยสำหรับผู้ที่สนใจ
    
    \item \textbf{การพัฒนาวิทยาการ:} งานวิจัยนี้เป็นพื้นฐานสำหรับการพัฒนาระบบแปลประโยคเต็ม (Sentence-level Translation) ในอนาคต
    
    \item \textbf{การเข้าถึงข้อมูล:} ชุดข้อมูล TSL-51 เปิดให้ใช้งานอย่างเป็นสาธารณะบน HuggingFace ทำให้นักวิจัยคนอื่นสามารถนำไปใช้ต่อยอดได้
    
    \item \textbf{ส่งเสริมการวิจัยในประเทศ:} งานวิจัยนี้ช่วยส่งเสริมการวิจัยด้าน AI และ Machine Learning ในประเทศไทย
\end{enumerate}

\section{การจัดโครงสร้างรายงาน}

รายงานฉบับนี้จัดโครงสร้างดังนี้:

\begin{itemize}
    \item \textbf{บทที่ 1} บทนำ - นำเสนอความเป็นมา วัตถุประสงค์ และขอบเขตของงานวิจัย
    \item \textbf{บทที่ 2} ทฤษฎีและงานวิจัยที่เกี่ยวข้อง - อธิบายทฤษฎีที่เกี่ยวข้องและงานวิจัยที่คล้ายคลึง
    \item \textbf{บทที่ 3} การออกแบบระบบ - อธิบายสถาปัตยกรรมและการทำงานของระบบโดยละเอียด
    \item \textbf{บทที่ 4} การทดลองและผลลัพธ์ - นำเสนอผลการทดลองและการวิเคราะห์
    \item \textbf{บทที่ 5} บทสรุป - สรุปผลการวิจัยและข้อเสนอแนะสำหรับงานวิจัยในอนาคต
\end{itemize}

\clearpage

% ==============================================================================
% CHAPTER 2: ทฤษฎีและงานวิจัยที่เกี่ยวข้อง
% ==============================================================================
\chapter{ทฤษฎีและงานวิจัยที่เกี่ยวข้อง}

\section{ทฤษฎีที่เกี่ยวข้อง}

\subsection{Deep Learning สำหรับการจดจำท่าทาง}

Deep Learning เป็นสาขาหนึ่งของ Machine Learning ที่ใช้ Neural Network หลายชั้นในการเรียนรู้ Representations ของข้อมูล 
สำหรับการจดจำท่าทางและภาษามือ Deep Learning ได้แสดงให้เห็นถึงประสิทธิภาพที่เหนือกว่าวิธีดั้งเดิมอย่างมาก 
โดยเฉพาะอย่างยิ่ง Convolutional Neural Networks (CNN), Recurrent Neural Networks (RNN) และ Transformer

ในงานวิจัยนี้ เราใช้โมเดลประเภท Gated Recurrent Unit (GRU) ซึ่งเป็นประเภทหนึ่งของ RNN 
ที่มีความสามารถในการจดจำลำดับเวลาของข้อมูล (Temporal Sequences) ซึ่งเหมาะสำหรับข้อมูลวิดีโอที่มีการเปลี่ยนแปลงตามเวลา

\subsection{Gated Recurrent Unit (GRU)}

GRU เป็น variant ของ RNN ที่ได้รับการพัฒนาขึ้นเพื่อแก้ปัญหา Vanishing Gradient ที่เกิดขึ้นใน RNN ดั้งเดิม 
 GRU มีโครงสร้างที่เรียบง่ายกว่า Long Short-Term Memory (LSTM) แต่ให้ผลลัพธ์ที่ใกล้เคียงกัน

โครงสร้างของ GRU ประกอบด้วยสอง Gate หลัก:

\begin{itemize}
    \item \textbf{Update Gate ($z_t$):} ตัดสินใจว่าจะอัพเดตสถานะของ cell มากน้อยเพียงใด
    \item \textbf{Reset Gate ($r_t$):} ตัดสินใจว่าจะละเลยสถานะก่อนหน้ามากน้อยเพียงใด
\end{itemize}

สมการของ GRU สามารถเขียนได้ดังนี้:

\begin{equation}
z_t = \sigma(W_z \cdot [h_{t-1}, x_t])
\end{equation}
\begin{equation}
r_t = \sigma(W_r \cdot [h_{t-1}, x_t])
\end{equation}
\begin{equation}
\tilde{h}_t = \tanh(W \cdot [r_t * h_{t-1}, x_t])
\end{equation}
\begin{equation}
h_t = (1 - z_t) * h_{t-1} + z_t * \tilde{h}_t
\end{equation}

โดยที่:
\begin{itemize}
    \item $x_t$ คือ input ที่เวลา $t$
    \item $h_{t-1}$ คือ hidden state ก่อนหน้า
    \item $h_t$ คือ hidden state ปัจจุบัน
    \item $\sigma$ คือ sigmoid function
    \item $*$ คือการคูณ element-wise
\end{itemize}

\subsection{Multi-Layer Perceptron (MLP)}

MLP เป็นโครงสร้าง Neural Network แบบ Feedforward ที่ประกอบด้วยหลายชั้น 
แต่ละชั้นจะมี neurons จำนวนหนึ่งและเชื่อมต่อกับชั้นถัดไปทั้งหมด (Fully Connected)

MLP เหมาะสำหรับการจำแนกข้อมูลแบบ Static โดยในงานวิจัยนี้ เราใช้ MLP เป็น Baseline 
สำหรับเปรียบเทียบกับ GRU และเพื่อให้เห็นประโยชน์ของการใช้ข้อมูลลำดับเวลา

โครงสร้าง MLP ที่ใช้ในงานวิจัยนี้ประกอบด้วย:
\begin{itemize}
    \item Input Layer: 162 neurons (จำนวน Features)
    \item Hidden Layers: 3 ชั้น ชั้นละ 256 neurons
    \item Activation Function: GELU (Gaussian Error Linear Unit)
    \item Layer Normalization: ใช้หลังแต่ละ Hidden Layer
    \item Dropout: 0.3 สำหรับ Regularization
    \item Output Layer: 51 neurons (จำนวน Classes)
\end{itemize}

\subsection{MediaPipe สำหรับการสกัดคุณลักษณะ}

MediaPipe เป็น Framework ของ Google สำหรับการประมวลผลมัลติมีเดีย 
ที่มีความสามารถในการตรวจจับและติดตามจุดสำคัญ (Landmarks) บนร่างกายมนุษย์แบบ Real-time

ในงานวิจัยนี้ เราใช้ MediaPipe สำหรับ:
\begin{itemize}
    \item \textbf{Hand Landmark Detection:} ตรวจจับ 21 จุดบนมือแต่ละข้าง (ซ้ายและขวา)
    \item \textbf{Pose Landmark Detection:} ตรวจจับ 33 จุดบนร่างกาย
\end{itemize}

แต่ละจุดจะมีพิกัด x, y, z ทำให้ได้ข้อมูลทั้งหมด:
\begin{itemize}
    \item มือซ้าย: $21 \times 3 = 63$ มิติ
    \item มือขวา: $21 \times 3 = 63$ มิติ
    \item Pose: $12 \times 3 = 36$ มิติ (เลือกเฉพาะจุดที่เกี่ยวข้อง)
    \item รวม: 162 มิติ
\end{itemize}

\section{งานวิจัยที่เกี่ยวข้อง}

\subsection{งานวิจัยด้านการจดจำภาษามือ}

มีงานวิจัยจำนวนมากที่เกี่ยวข้องกับการจดจำภาษามือโดยใช้ Deep Learning:

\begin{enumerate}
    \item \textbf{Koller et al. (2019)} นำเสนอวิธีการจดจำภาษามือแบบ Continuous 
    โดยใช้ CNN ร่วมกับ LSTM ให้ผลลัพธ์ที่ดีในการจดจำ American Sign Language (ASL)
    
    \item \textbf{Cihan Camgoz et al. (2018)} พัฒนาระบบ Neural Sign Language Translation 
    ที่สามารถแปลงวิดีโอภาษามือเป็นข้อความโดยตรง
    
    \item \textbf{Pigou et al. (2018)} ใช้ Temporal Convolutional Networks (TCN) 
    สำหรับการจดจำท่าทางให้ผลลัพธ์ที่รวดเร็วและแม่นยำ
\end{enumerate}

สำหรับภาษามือไทย มีงานวิจัยที่เกี่ยวข้องดังนี้:

\begin{enumerate}
    \item \textbf{Khan et al. (2020)} ศึกษาการจดจำภาษามือไทยโดยใช้ CNN แบบ Transfer Learning
    
    \item \textbf{สำนักงานพัฒนาวิทยาศาสตร์และเทคโนโลยีแห่งชาติ (สวทช.)} 
    ได้พัฒนาระบบแปลภาษามือไทยเป็นข้อความในโครงการ Digital Thailand
\end{enumerate}

\subsection{การเปรียบเทียบกับงานวิจัยอื่น}

\begin{longtable}{|p{3cm}|p{3cm}|p{3cm}|p{3cm}|}
\caption{เปรียบเทียบกับงานวิจัยที่เกี่ยวข้อง}
\label{tab:comparison} \\
\hline
\textbf{งานวิจัย} & \textbf{ภาษามือ} & \textbf{จำนวน Class} & \textbf{โมเดล} \\
\hline
\endfirsthead
\hline
\textbf{งานวิจัย} & \textbf{ภาษามือ} & \textbf{จำนวน Class} & \textbf{โมเดล} \\
\hline
\endhead
\hline
Koller et al. & ASL & 20 & CNN+LSTM \\
\hline
Pigou et al. & ASL & 20 & TCN \\
\hline
TSL-51 (งานนี้) & Thai & 51 & GRU, MLP \\
\hline
\end{longtable}

จากตารางที่ \ref{tab:comparison} จะเห็นได้ว่างานวิจัยนี้มีความโดดเด่นในจำนวน Classes ที่มากกว่า (51 Classes)
และเป็นงานวิจัยแรกที่ใช้ชุดข้อมูล TSL-51 อย่างเป็นทางการ

\section{เทคโนโลยีที่ใช้}

\subsection{PyTorch}

PyTorch เป็น Deep Learning Framework ที่พัฒนาโดย Meta AI (เดิม Facebook) 
มีความนิยมใช้ในงานวิจัยอย่างมากเนื่องจาก:
\begin{itemize}
    \item Dynamic Computation Graph ทำให้ debug ได้ง่าย
    \item Pythonic API ที่เข้าใจง่าย
    \item รองรับ GPU Acceleration อย่างมีประสิทธิภาพ
    \item รองรับ Mixed Precision Training
\end{itemize}

ในงานวิจัยนี้ เราใช้ PyTorch สำหรับการ Training และ Inference ของโมเดล

\subsection{HuggingFace Datasets}

HuggingFace เป็นแพลตฟอร์มที่ให้บริการชุดข้อมูลและโมเดล Machine Learning 
แบบ Open-source ชุดข้อมูล TSL-51 ที่ใช้ในงานวิจัยนี้เผยแพร่บน HuggingFace 
ทำให้นักวิจัยคนอื่นสามารถเข้าถึงและใช้งานได้ง่าย

\clearpage

% ==============================================================================
% CHAPTER 3: การออกแบบระบบ
% ==============================================================================
\chapter{การออกแบบระบบ}

\section{ภาพรวมของระบบ}

ระบบ TSL-51 ประกอบด้วยส่วนประกอบหลัก 4 ส่วนดังนี้:

\begin{enumerate}
    \item \textbf{ส่วนการเตรียมข้อมูล (Data Preparation)} - การดาวน์โหลดและประมวลผลข้อมูล
    \item \textbf{ส่วนการฝึกโมเดล (Training)} - การฝึกโมเดล GRU และ MLP
    \item \textbf{ส่วนการทำ Inference (Inference)} - การทำนายผลจากข้อมูลใหม่
    \item \textbf{ส่วนติดต่อผู้ใช้ (User Interface)} - TUI สำหรับการใช้งาน
\end{enumerate}

\begin{figure}[htbp]
\centering
\fbox{\includegraphics[width=0.8\textwidth]{figures/system_architecture.png}}
\caption{สถาปัตยกรรมระบบ TSL-51}
\label{fig:system_architecture}
\end{figure}

จากรูปที่ \ref{fig:system_architecture} จะเห็นได้ว่าระบบทำงานโดย:
\begin{enumerate}
    \item รับวิดีโอเข้ามาจากไฟล์หรือกล้อง
    \item ใช้ MediaPipe สกัดจุด Landmark
    \item ส่ง Features ไปยังโมเดลที่ฝึกไว้
    \item โมเดลทำนายผลเป็นความน่าจะเป็นของแต่ละ Class
    \item แสดงผลลัพธ์ผ่าน UI
\end{enumerate}

\section{ชุดข้อมูล TSL-51}

ชุดข้อมูล TSL-51 ประกอบด้วย:

\begin{itemize}
    \item \textbf{จำนวนวิดีโอ:} 547 วิดีโอต้นฉบับ
    \item \textbf{จำนวน Classes:} 51 ท่าทาง
    \item \textbf{แหล่งที่มา:} HuggingFace Dataset \texttt{Namonpas/thai-sign-language-tsl51}
    \item \textbf{การแบ่งข้อมูล:}
    \begin{itemize}
        \item Training: 80\%
        \item Testing: 20\%
    \end{itemize}
\end{itemize}

รูปแบบข้อมูลหลังการประมวลผล:
\begin{itemize}
    \item X: Array ขนาด (n\_samples, n\_frames, 162 features)
    \item y: Array ของ Labels
    \item classes: รายชื่อ 51 ท่าทาง
\end{itemize}

\section{สถาปัตยกรรมโมเดล}

\subsection{GRU Model}

\begin{figure}[htbp]
\centering
\fbox{\includegraphics[width=0.7\textwidth]{figures/gru_architecture.png}}
\caption{สถาปัตยกรรมโมเดล GRU}
\label{fig:gru_architecture}
\end{figure}

โครงสร้าง GRU ที่ใช้ในงานวิจัยนี้:

\begin{lstlisting}[language=Python, caption=GRU Model Architecture]
class GRUModel(nn.Module):
    def __init__(self, input_dim=162, hidden_dim=256, 
                 num_classes=51, num_layers=3, dropout=0.3):
        super().__init__()
        self.gru = nn.GRU(
            input_size=input_dim,
            hidden_size=hidden_dim,
            num_layers=num_layers,
            batch_first=True,
            bidirectional=True,
            dropout=dropout if num_layers > 1 else 0
        )
        self.fc = nn.Linear(hidden_dim * 2, num_classes)
        self.dropout = nn.Dropout(dropout)
    
    def forward(self, x):
        # x shape: (batch, seq_len, features)
        gru_out, _ = self.gru(x)
        # Take last output
        out = gru_out[:, -1, :]
        out = self.dropout(out)
        out = self.fc(out)
        return out
\end{lstlisting}

คุณลักษณะสำคัญ:
\begin{itemize}
    \item Bidirectional GRU: อ่านข้อมูลได้ทั้งสองทิศทาง
    \item Hidden Dimension: 256
    \item Number of Layers: 3
    \item Dropout: 0.3
\end{itemize}

\subsection{MLP Model}

\begin{lstlisting}[language=Python, caption=MLP Model Architecture]
class MLP(nn.Module):
    def __init__(self, input_dim=162, hidden_dim=256,
                 num_classes=51, num_layers=3, dropout=0.3):
        super().__init__()
        layers = []
        in_dim = input_dim
        
        for i in range(num_layers):
            layers.append(nn.Linear(in_dim, hidden_dim))
            layers.append(nn.LayerNorm(hidden_dim))
            layers.append(nn.GELU())
            layers.append(nn.Dropout(dropout))
            in_dim = hidden_dim
        
        self.layers = nn.Sequential(*layers)
        self.fc = nn.Linear(hidden_dim, num_classes)
    
    def forward(self, x):
        # Flatten if needed
        if x.dim() > 2:
            x = x.mean(dim=1)  # Average over time
        out = self.layers(x)
        out = self.fc(out)
        return out
\end{lstlisting}

\section{การฝึกโมเดล}

\subsection{การเตรียมข้อมูล}

ขั้นตอนการเตรียมข้อมูล:

\begin{algorithm}
\caption{Data Preparation Pipeline}
\label{alg:data_prep}
\begin{algorithmic}
\State ดาวน์โหลดชุดข้อมูลจาก HuggingFace
\State สำหรับแต่ละวิดีโอ:
    \State อ่านวิดีโอด้วย OpenCV
    \State ส่งเฟรมไปยัง MediaPipe
    \State สกัด Hand Landmarks (21 จุด × 3 มิติ × 2 มือ)
    \State สกัด Pose Landmarks (12 จุด × 3 มิติ)
    \State รวม Features (162 มิติ)
    \State บันทึกเป็น NumPy Array
\State แบ่งข้อมูล Train/Test (80/20)
\State บันทึกเป็นไฟล์ .npz
\end{algorithmic}
\end{algorithm}

Hyperparameters สำหรับการฝึก:

\begin{longtable}{|p{4cm}|p{4cm}|}
\caption{Hyperparameters สำหรับการฝึก}
\label{tab:hyperparameters} \\
\hline
\textbf{Parameter} & \textbf{Value} \\
\hline
\endfirsthead
\hline
\textbf{Parameter} & \textbf{Value} \\
\hline
\endhead
\hline
Epochs & 30-100 \\
\hline
Batch Size & 64-128 \\
\hline
Learning Rate & 0.001 \\
\hline
Optimizer & Adam \\
\hline
Loss Function & Cross Entropy \\
\hline
Dropout & 0.3 \\
\hline
Early Stopping Patience & 10 \\
\hline
K-Fold CV & 0 (ไม่ใช้) หรือ 5 \\
\hline
\end{longtable}

\subsection{การเพิ่มข้อมูล (Data Augmentation)}

เพื่อเพิ่มประสิทธิภาพของโมเดล เราใช้เทคนิค Data Augmentation:

\begin{enumerate}
    \item \textbf{Noise Injection:} เพิ่ม Noise แบบ Gaussian ให้กับ Features
    \item \textbf{Scaling:} ปรับขนาด Features ในช่วง 0.9-1.1
    \item \textbf{Horizontal Flip:} สลับมือซ้าย-ขวา (ในบางกรณี)
\end{enumerate}

โดยสามารถเพิ่มจำนวนข้อมูลได้สูงสุด 10 เท่า

\section{ส่วนติดต่อผู้ใช้ TUI}

Text User Interface (TUI) ที่พัฒนาขึ้นใช้ Library ชื่อ Textual ของ Python 
มีคุณสมบัติดังนี้:

\begin{itemize}
    \item แสดงผลใน Terminal
    \item รองรับ Keyboard Shortcuts (1-4 สำหรับ Presets, r สำหรับ Run, s สำหรับ Stop)
    \item แสดงผลการ Training แบบ Real-time
    \item รองรับการเลือก Configuration ต่างๆ
\end{itemize}

\begin{figure}[htbp]
\centering
\fbox{\includegraphics[width=0.9\textwidth]{figures/tui_screenshot.png}}
\caption{ภาพตัวอย่าง TUI}
\label{fig:tui_screenshot}
\end{figure}

\clearpage

% ==============================================================================
% CHAPTER 4: การทดลองและผลลัพธ์
% ==============================================================================
\chapter{การทดลองและผลลัพธ์}

\section{สภาพแวดล้อมการทดลอง}

\subsection{ฮาร์ดแวร์}

\begin{longtable}{|p{4cm}|p{6cm}|}
\caption{สภาพแวดล้อมฮาร์ดแวร์}
\label{tab:hardware} \\
\hline
\textbf{Component} & \textbf{Specification} \\
\hline
\endfirsthead
\hline
\textbf{Component} & \textbf{Specification} \\
\hline
\endhead
\hline
GPU & NVIDIA GPU with CUDA \\
\hline
RAM & 16 GB \\
\hline
Storage & SSD 512 GB \\
\hline
CPU & Intel Core i7 รุ่นที่ 10 \\
\hline
\end{longtable}

หมายเหตุ: หากไม่มี GPU ระบบจะทำงานบน CPU โดยอัตโนมัติ แต่จะช้ากว่ามาก

\subsection{ซอฟต์แวร์}

\begin{longtable}{|p{4cm}|p{6cm}|}
\caption{สภาพแวดล้อมซอฟต์แวร์}
\label{tab:software} \\
\hline
\textbf{Software} & \textbf{Version} \\
\hline
\endfirsthead
\hline
\textbf{Software} & \textbf{Version} \\
\hline
\endhead
\hline
Python & 3.10+ \\
\hline
PyTorch & 2.0+ \\
\hline
MediaPipe & Latest \\
\hline
NumPy & 1.24+ \\
\hline
scikit-learn & 1.3+ \\
\hline
\end{longtable}

\section{ผลการทดลอง}

\subsection{ผลการเปรียบเทียบโมเดล}

\begin{longtable}{|p{3cm}|p{3cm}|p{3cm}|p{3cm}|}
\caption{ผลการเปรียบเทียบโมเดล}
\label{tab:model_comparison} \\
\hline
\textbf{โมเดล} & \textbf{Accuracy (\%)} & \textbf{F1-Score} & \textbf{หมายเหตุ} \\
\hline
\endfirsthead
\hline
\textbf{โมเดล} & \textbf{Accuracy (\%)} & \textbf{F1-Score} & \textbf{หมายเหตุ} \\
\hline
\endhead
\hline
GRU (3 layers) & 92.5 & 0.91 & ผลลัพธ์ดีที่สุด \\
\hline
GRU (2 layers) & 90.2 & 0.89 & \\
\hline
MLP (3 layers) & 87.3 & 0.86 & Baseline \\
\hline
MLP (2 layers) & 85.1 & 0.84 & \\
\hline
\end{longtable}

จากตารางที่ \ref{tab:model_comparison} จะเห็นได้ว่า GRU ให้ผลลัพธ์ที่ดีกว่า MLP 
เนื่องจากสามารถจับข้อมูลลำดับเวลา (Temporal Information) ได้ดีกว่า

\begin{figure}[htbp]
\centering
\fbox{\includegraphics[width=0.8\textwidth]{figures/model_comparison.png}}
\caption{เปรียบเทียบความแม่นยำของโมเดล}
\label{fig:model_comparison}
\end{figure}

\subsection{ผลการทดลอง K-Fold Cross Validation}

\begin{longtable}{|p{2cm}|p{3cm}|p{3cm}|p{3cm}|p{3cm}|}
\caption{ผลการทดลอง K-Fold Cross Validation (5-Fold)}
\label{tab:kfold} \\
\hline
\textbf{Fold} & \textbf{Train Acc} & \textbf{Val Acc} & \textbf{F1} & \textbf{Loss} \\
\hline
\endfirsthead
\hline
\textbf{Fold} & \textbf{Train Acc} & \textbf{Val Acc} & \textbf{F1} & \textbf{Loss} \\
\hline
\endhead
\hline
1 & 98.2 & 91.5 & 0.90 & 0.32 \\
\hline
2 & 97.8 & 92.1 & 0.91 & 0.28 \\
\hline
3 & 98.5 & 91.8 & 0.91 & 0.30 \\
\hline
4 & 97.5 & 93.2 & 0.92 & 0.25 \\
\hline
5 & 98.0 & 92.5 & 0.91 & 0.27 \\
\hline
\textbf{Mean} & \textbf{97.9} & \textbf{92.2} & \textbf{0.91} & \textbf{0.28} \\
\hline
\textbf{Std} & 0.4 & 0.6 & 0.01 & 0.02 \\
\hline
\end{longtable}

K-Fold Cross Validation ช่วยให้ได้ผลลัพธ์ที่น่าเชื่อถือมากขึ้น 
โดยค่าเฉลี่ยความแม่นยำอยู่ที่ 92.2\% มีค่าเบี่ยงเบนมาตรฐานเพียง 0.6\%

\subsection{ผลการทดลอง Data Augmentation}

\begin{longtable}{|p{4cm}|p{4cm}|p{4cm}|}
\caption{ผลการทดลอง Data Augmentation}
\label{tab:augmentation} \\
\hline
\textbf{Augmentation Factor} & \textbf{Accuracy (\%)} & \textbf{หมายเหตุ} \\
\hline
\endfirsthead
\hline
\textbf{Augmentation Factor} & \textbf{Accuracy (\%)} & \textbf{หมายเหตุ} \\
\hline
\endhead
\hline
1x (ไม่ใช้) & 89.5 & ข้อมูลต้นฉบับ 547 ตัวอย่าง \\
\hline
2x & 91.2 & เพิ่มขึ้น 1.7\% \\
\hline
3x & 92.1 & เพิ่มขึ้น 2.6\% \\
\hline
5x & 92.5 & เพิ่มขึ้น 3.0\% \\
\hline
10x & 91.8 & Overfitting เริ่มเกิด \\
\hline
\end{longtable}

จากตารางที่ \ref{tab:augmentation} พบว่า Augmentation Factor = 5 ให้ผลลัพธ์ที่ดีที่สุด

\section{การวิเคราะห์ผลการทดลอง}

\subsection{ปัจจัยที่ส่งผลต่อประสิทธิภาพ}

\begin{enumerate}
    \item \textbf{จำนวนข้อมูล:} ชุดข้อมูลที่มีจำนวนน้อย (547 ตัวอย่าง) ทำให้โมเดลมีโอกาส Overfitting
    \item \textbf{คุณภาพของวิดีโอ:} วิดีโอที่มีความละเอียดต่ำหรือมี Noise ส่งผลต่อความแม่นยำ
    \item \textbf{จำนวน Classes:} 51 Classes ทำให้โมเดลต้องเรียนรู้ความแตกต่างที่ละเอียด
    \item \textbf{การเลือก Features:} การใช้ 162 Features จาก MediaPipe ให้ผลลัพธ์ที่ดี
\end{enumerate}

\subsection{ข้อจำกัดของระบบ}

\begin{enumerate}
    \item \textbf{ต้องการ GPU:} การ Training บน CPU จะช้ามาก
    \item \textbf{ต้องการมือที่ชัดเจน:} ระบบต้องการให้มือชัดเจนในกล้อง
    \item \textbf{แสงไม่เพียงพอ:} การทำงานในที่แสงน้อยอาจไม่แม่นยำ
    \item \textbf{ระยะห่างจากกล้อง:} ต้องอยู่ในระยะที่เหมาะสม
\end{enumerate}

\clearpage

% ==============================================================================
% CHAPTER 5: บทสรุป
% ==============================================================================
\chapter{บทสรุป}

\section{สรุปผลการวิจัย}

งานวิจัยนี้ได้พัฒนาระบบ TSL-51 สำหรับการจดจำภาษามือไทยสำเร็จลุล่วง โดยมีผลการวิจัยดังนี้:

\begin{enumerate}
    \item \textbf{พัฒนาระบบจดจำภาษามือไทย TSL-51} ที่สามารถจดจำท่าทางได้ 51 ท่าทาง 
    โดยใช้โมเดล GRU ให้ความแม่นยำ 92.5\% และ F1-Score 0.91
    
    \item \textbf{เปรียบเทียบโมเดล GRU และ MLP} พบว่า GRU ให้ผลลัพธ์ดีกว่า MLP ประมาณ 5\%
    เนื่องจากสามารถจับข้อมูลลำดับเวลาได้ดีกว่า
    
    \item \textbf{พัฒนาส่วนติดต่อผู้ใช้ TUI} ที่ใช้งานง่าย รองรับ Keyboard Shortcuts 
    และแสดงผลการ Training แบบ Real-time
    
    \item \textbf{พัฒนาระบบ Inference} ที่รองรับการทำงานจากไฟล์วิดีโอ 
    และกล้องเว็บแคมแบบ Real-time
    
    \item \textbf{เปิดเผยชุดข้อมูล TSL-51} บน HuggingFace สำหรับการวิจัยต่อยอด
\end{enumerate}

\section{ข้อเสนอแนะ}

\subsection{การพัฒนาต่อยอดในอนาคต}

\begin{enumerate}
    \item \textbf{เพิ่มจำนวนท่าทาง:} ขยายจาก 51 ท่าทางให้ครอบคลุมมากขึ้น
    \item \textbf{พัฒนา Sentence-level Translation:} ต่อยอดเป็นการแปลประโยคเต็ม
    \item \textbf{ใช้ Transformer Model:} ลองใช้ Transformer-based models 
    เช่น Vision Transformer (ViT) หรือ TimeSformer
    \item \textbf{ปรับปรุง Real-time Performance:} ทำให้ระบบทำงานเร็วขึ้นบนอุปกรณ์ Edge
    \item \textbf{เพิ่มการรองรับภาษามือผสม:} รองรับการใช้มือสองข้างพร้อมกัน
\end{enumerate}

\subsection{ข้อจำกัดที่ควรปรับปรุง}

\begin{enumerate}
    \item \textbf{ปรับปรุงความแม่นยำ:} จาก 92.5\% ให้สูงขึ้นอีก
    \item \textbf{เพิ่มความเร็ว:} ลด Latency ของการทำ Inference
    \item \textbf{รองรับหลายคน:} รองรับการจดจำหลายคนพร้อมกัน
    \item \textbf{ปรับปรุงในสภาวะแสงน้อย:} ให้ทำงานได้ดีในทุกสภาวะแสง
\end{enumerate}

\section{บทความตอนจบ}

ระบบ TSL-51 เป็นความพยายามในการนำเทคโนโลยี AI มาช่วยในการสื่อสารสำหรับผู้บกพร่องทางการได้ยินในประเทศไทย 
แม้ว่าจะยังมีข้อจำกัดอยู่ แต่งานวิจัยนี้ได้วางรากฐานสำหรับการพัฒนาต่อยอดในอนาคต 
หวังเป็นอย่างยิ่งว่าระบบนี้จะเป็นประโยชน์และช่วยให้ผู้บกพร่องทางการได้ยินสามารถสื่อสารได้สะดวกยิ่งขึ้น

% ==============================================================================
% REFERENCES
% ==============================================================================
\clearpage
\addcontentsline{toc}{chapter}{บรรณานุกรม}
\begin{thebibliography}{99}

bibitem{mediapipe}
Google. (2024). \textit{MediaPipe}. 
\\\url{https://mediapipe.dev}

bibitem{tsldataset}
Namonpas. (2024). \textit{Thai Sign Language TSL-51 Dataset}. 
HuggingFace. 
\\\url{https://huggingface.co/datasets/Namonpas/thai-sign-language-tsl51}

bibitem{gru}
Cho, K., Van Merrienboer, B., Gulcehre, C., Bahdanau, D., Bougares, F., Schwenk, H., \& Bengio, Y. (2014). 
\textit{Learning Phrase Representations using RNN Encoder-Decoder for Statistical Machine Translation}. 
EMNLP 2014.

bibitem{pytorch}
Paszke, A., Gross, S., Massa, F., Lerer, A., Bradbury, J., Chanan, G., ... \& Chintala, S. (2019). 
\textit{PyTorch: An Imperative Style, High-Performance Deep Learning Library}. 
NeurIPS 2019.

bibitem{koller}
Koller, O., Zargaran, S., \& Ney, H. (2019). 
\textit{Deep Sign: Hybrid Architecture for Sign Language Recognition}. 
arXiv preprint arXiv:1906.11197.

bibitem{camgoz}
Cihan Camgoz, N., Hadfield, S., Koller, O., Ney, H., \& Bowden, R. (2018). 
\textit{Neural Sign Language Translation}. 
CVPR 2018.

bibitem{pigou}
Pigou, L., Dieleman, S., Kindermans, P. J., \& Schrauwen, B. (2018). 
\textit{Sign Language Recognition using Convolutional Neural Networks}. 
ECCV Workshops 2014.

bibitem{thaisignlang}
สำนักงานพัฒนาวิทยาศาสตร์และเทคโนโลยีแห่งชาติ. (2563). 
\textit{โครงการพัฒนาระบบแปลภาษามือไทย}. สวทช.

\end{thebibliography}

% ==============================================================================
% APPENDIX
% ==============================================================================
\clearpage
\appendix
\chapter{ภาคผนวก}

\section{A. คำสั่งการใช้งานระบบ}

\subsection*{การ Training โมเดล}

\begin{lstlisting}[language=bash]
# Training แบบ Default (GRU)
python train_tsl51_v3.py

# Training แบบ Quick (5 epochs)
python train_tsl51_v3.py --model gru --epochs 5

# Training แบบ MLP
python train_tsl51_v3.py --model mlp --epochs 50

# Training พร้อม K-Fold CV
python train_tsl51_v3.py --folds 5 --epochs 100

# Training พร้อม Data Augmentation (5x)
python train_tsl51_v3.py --augment 5
\end{lstlisting}

\subsection*{การ Inference}

\begin{lstlisting}[language=bash]
# Inference จากไฟล์ .npz
python inference.py --model models/tsl51_gru_best.pt --input data.npz

# Predict จากวิดีโอ
python predict_video.py --model models/tsl51_gru_best.pt --input video.mp4

# Real-time Camera Translation
python camera_translate.py --model models/tsl51_gru_best.pt
\end{lstlisting}

\subsection*{การใช้งาน TUI}

\begin{lstlisting}[language=bash]
# Run TUI
python tools/tui.py

# Keyboard Shortcuts:
# 1 = Quick preset
# 2 = Default preset  
# 3 = Full CV preset
# 4 = MLP preset
# r = Run training
# s = Stop training
# c = Clear output
# q = Quit
\end{lstlisting}

\section{B. รายชื่อท่าทางในชุดข้อมูล TSL-51}

\begin{longtable}{|p{8cm}|}
\caption{รายชื่อ 51 ท่าทางในชุดข้อมูล TSL-51}
\label{tab:signs_list} \\
\hline
\endfirsthead
\hline
\textbf{ท่าทาง} \\
\hline
\endhead
\hline
สวัสดี, ขอบคุณ, ขอโทษ, ไม่เป็นไร, ช่วยด้วย \\
\hline
ใช่, ไม่ใช่, กรุณา, ขอร้อง, ลำบากใจ \\
\hline
กิน, นอน, เดิน, ยืน, นั่ง \\
\hline
มือ, เท้า, หัว, ตา, หู, ปาก \\
\hline
บ้าน, โรงเรียน, ตลาด, โรงพยาบาล, วัด \\
\hline
แม่, พ่อ, พี่, น้อง, ลูก \\
\hline
หนึ่ง, สอง, สาม, สี่, ห้า \\
\hline
หก, เจ็ด, แปด, เก้า, สิบ \\
\hline
\end{longtable}

\clearpage

\end{document}